\exercisesfor{7}

第 7 节的约定除非另有说明均有效。

\begin{enumerate}
\def\labelenumi{\arabic{enumi}.}
\tightlist
\item
  \begin{enumerate}
  \def\labelenumii{(\alph{enumii})}
  \tightlist
  \item
    李 K-代数幂零的充分必要条件是它同构于某个代数 \(n(n, \mathbf{K})\) 的一个子代数。
  \end{enumerate}
\end{enumerate}

\begin{enumerate}
\def\labelenumi{(\alph{enumi})}
\setcounter{enumi}{1}
\tightlist
\item
  假设 K 代数闭。李 K-代数可解的充分必要条件是它同构于某个代数 \(t(n, \mathbf{K})\) 的一个子代数。
\end{enumerate}

\begin{enumerate}
\def\labelenumi{\arabic{enumi}.}
\setcounter{enumi}{1}
\item
  设 g 为李代数，n 为其最大幂零理想。则存在有限维向量空间 V 以及 g 到 \(\mathfrak{sl}(V)\) 的一个子代数上的同构，它把 n 的每个元素映为 V 的幂零自同态。（利用 Ado 定理与第 1 节练习 5。）
\item
  设 g 为李代数，U 为其包络代数。
\end{enumerate}

\begin{enumerate}
\def\labelenumi{(\alph{enumi})}
\item
  证明对一切 \(u \in \mathrm{U}\)（ \(u \neq 0\) ），存在一个有限维表示 \(\pi\)（\(\mathrm{U}\) 上的），使 \(\pi(u) \neq 0\)。（利用练习 2 与第 3 节练习 5 (b)。）
\item
  由 (a) 推出：若 g 半单，则存在无穷多个互不等价的 g 的有限维单表示。（设 \(\rho_{1},\ldots,\rho_{n}\) 为 g 的有限维单表示，\(N_{1},\ldots,N_{n}\) 为相应 U-模的零化理想，\(N=\bigcap_{i=1}^{n}N_{i}\)。证明 N 在 U 中有余限维。设 \(n\in N,n\neq0\)。对 n 应用 (a)。）
\end{enumerate}

\begin{enumerate}
\def\labelenumi{\arabic{enumi}.}
\setcounter{enumi}{3}
\tightlist
\item
  设 g 为李代数，a 为 g 的次不变子代数，\(\mathfrak{r}(a)\) 为 a 的根基，\(\rho\) 为 a 的有限维表示。\(\rho\) 容许到 g 的有限维扩张的充分必要条件是 \(\mathcal{D}_{\mathfrak{g}} \cap \mathfrak{r}(a)\) 包含在 \(\rho\) 的最大幂零性理想中。（为证充分性，如下进行：设 \(g = g_{0} \supset g_{1} \supset \cdots \supset g_{k} = a\) 为 g 的具有第 6 节练习 11 (b) 中诸性质的合成列。用归纳法证明存在扩张 \(\rho_{i}\)（\(\rho\) 到 \(g_{0}\) 的扩张），其最大幂零性理想包含 \(\mathcal{D}_{\mathfrak{g}} \cap \mathfrak{r}(g_{i})\)。在第 6 节练习 11 (b) 的记号下，由 \(\rho_{i+1}\) 到 \(\rho_{i}\) 的过渡由定理 1 得出：当 \(b_{i}\) 单，或当 \(b_{i}\) 为 1 维且
\end{enumerate}

\[
\mathscr {D} \mathfrak {g} \cap \mathfrak {r} (\mathfrak {g} _ {i + 1}) = \mathscr {D} \mathfrak {g} \cap \mathfrak {r} (\mathfrak {g} _ {i}).
\]

时成立。当 \(\mathcal{D}\mathfrak{g} \cap \mathfrak{r}(\mathfrak{g}_{i+1}) \neq \mathcal{D}\mathfrak{g} \cap \mathfrak{r}(\mathfrak{g}_i)\) 时，可设 \(\mathfrak{h}_i \subset \mathcal{D}\mathfrak{g} \cap \mathfrak{r}(\mathfrak{g}_i)\)。于是 \(\mathfrak{r}(\mathfrak{g}_i) = \mathfrak{r}(\mathfrak{g}) \cap \mathfrak{g}_i\)（第 5 节练习 10），从而可应用定理 1。）

¶ 5. 设 K 为特征 p \textgreater{} 0 的域，g 为 K 上有限维 n 李代数，U 为 g 的包络代数。

\begin{enumerate}
\def\labelenumi{(\alph{enumi})}
\item
  \(p\) -多项式（\(\mathbf{K}\) 上、一个不定元的）是指 \(\mathrm{K}[\mathrm{X}]\) 中仅有的 \(\neq 0\) 项都是 \(p\) 的方幂次数的多项式。 证明对每个多项式 \(f(\mathbf{X}) \neq 0\)（在 \(\mathrm{K}[\mathrm{X}]\) 中），存在 \(g(\mathbf{X}) \in \mathrm{K}[\mathbf{X}]\) ，使 \(f(\mathbf{X})g(\mathbf{X})\) 是 \(p\) -多项式（考察诸 \(\mathbf{X}^{p_k}\) 被 \(f(\mathbf{X})\) 作欧几里得除法所得的余式）。
\item
  证明对每个元素 \(z \in \mathfrak{g}\)，存在非零多项式 \(f(\mathbf{X}) \in \mathbf{K}[\mathbf{X}]\)，使 \(f(z)\) 属于 U 的中心 C。（考察自同态 ad \(z\) 的最小多项式，对此多项式应用 (a)，并利用第 1 节练习 19 的公式 (1)。）
\item
  设 \((e_i)_{1\leqslant i\leqslant n}\) 为 \(\mathfrak{g}\) 的基。对每个 \(i\)，设 \(f_{i}\neq 0\) 为 \(\mathrm{K}[\mathrm{X}]\) 中使 \(f_{i}(e_{i})\in \mathrm{C}\) 的多项式；设 \(d_{i}\) 为其次数，\(\mathrm{L}\) 为 \(\mathrm{U}\) 中由诸 \(y_{i} = f_{i}(e_{i})\) 生成的双边理想。证明模 \(\mathrm{L}\) 的、元素 \(e_1^{\alpha_1}\ldots e_n^{\alpha_n}\)（其中 \(0\leqslant \alpha_{i} < d_{i}\)）的类构成 \(\mathrm{U / L}\) 的基。
\item
  证明 \(\mathfrak{g}\) 容许有限维忠实的线性表示（选取诸 \(f_{i}\)，使 \(d_{i} > 1\) 对一切 \(i\) 成立）。
\item
  证明若 \(g \neq \{0\}\)，则存在 g 的非半单有限维线性表示。（设诸 \(f_{t}\) 的次数为 \(d_{t} > 1\)，在 L 的定义中把诸 \(f_{t}\) 换成 \(f_{t}^{2}\)，并注意此时 U/L 的中心中存在 \(\neq 0\) 的幂零元素，因而 U/L 不是半单的。）
\end{enumerate}


